% answers.tex —— 2014 年·四川卷·文科 · 参考答案与详细解析片段
% 由 tools/generate_projects.py 依来源总账自 2014/2014_数学真题_四川文_试题.md + 2014/2014_数学真题_四川文_解析正文.md + 2014/2014_数学真题_四川文_答案要点.md 生成，勿手改（改源 MD 后重跑）


\begin{answer}
\gksolhead{1}{}

\textbf{题目}：

已知集合 $A=\{x\mid (x+1)(x-2)\leqslant 0\}$，集合 $B$ 为整数集，则 $A\cap B=$\mbox{（　　）}

（A）$\{-1,0\}$；（B）$\{0,1\}$；（C）$\{-2,-1,0,1\}$；（D）$\{-1,0,1,2\}$。

\Needspace{4\baselineskip}
\textbf{解析}

\Needspace{5\baselineskip}
由 $(x+1)(x-2)\leqslant 0$ 得
\nopagebreak[4]
\[-1\leqslant x\leqslant 2,\]
\Needspace{5\baselineskip}
即 $A=[-1,2]$。区间内的整数为 $-1,0,1,2$，故
\nopagebreak[4]
\[A\cap B=\{-1,0,1,2\}.\]

选 D。

\gkanswer{答案：D。$(x+1)(x-2)\leqslant 0$ 给出 $-1\leqslant x\leqslant 2$，其中的整数为 $-1,0,1,2$。}

\end{answer}
\begin{answer}
\gksolhead{2}{}

\textbf{题目}：

在“世界读书日”前夕，为了了解某地 $5000$ 名居民某天的阅读时间，从中抽取了 $200$ 名居民的阅读时间进行统计分析。在这个问题中，$5000$ 名居民的阅读时间的全体是\mbox{（　　）}

（A）总体；（B）个体；（C）样本的容量；（D）从总体中抽取的一个样本。

\Needspace{4\baselineskip}
\textbf{解析}

考察对象是某地 $5000$ 名居民某天的阅读时间，这 $5000$ 个数据的全体构成总体；从中抽取的 $200$ 个数据是样本，$200$ 是样本容量，每一个居民的阅读时间是个体。题干问“$5000$ 名居民的阅读时间的全体”，即总体。

选 A。

\gkanswer{答案：A。$5000$ 名居民阅读时间的全体是考察对象的全体，即总体。}

\end{answer}
\begin{answer}
\gksolhead{3}{}

\textbf{题目}：

为了得到函数 $y=\sin(x+1)$ 的图象，只需把函数 $y=\sin x$ 的图象上所有的点\mbox{（　　）}

（A）向左平行移动 $1$ 个单位长度；（B）向右平行移动 $1$ 个单位长度；（C）向左平行移动 $\pi$ 个单位长度；（D）向右平行移动 $\pi$ 个单位长度。

\Needspace{4\baselineskip}
\textbf{解析}

\Needspace{5\baselineskip}
把 $y=\sin x$ 的图象上所有点向左平移 $1$ 个单位长度，所得图象的解析式为
\nopagebreak[4]
\[y=\sin\bigl(x+1\bigr),\]
恰为所求。一般地，$y=f(x)$ 向左平移 $a>0$ 个单位得到 $y=f(x+a)$。

选 A。

\gkanswer{答案：A。$y=\sin(x+1)=\sin\bigl(x-(-1)\bigr)$，故只需向左平移 $1$ 个单位长度。}

\end{answer}
\begin{answer}
\gksolhead{4}{}

\textbf{题目}：

某三棱锥的侧视图、俯视图如图所示，则该三棱锥的体积是\mbox{（　　）}

（锥体体积公式：$V=\dfrac13Sh$，其中 $S$ 为底面面积，$h$ 为高）

（A）$3$；（B）$2$；（C）$\sqrt3$；（D）$1$。
\nopagebreak[4]
\begin{center}% ============================================================
% f2014_sichuan_arts_q04_three_views.tex
% 2014 四川卷（文科）第 4 题：三棱锥的侧视图与俯视图
% 来源 PDF：普通高考/2014/2014四川文.pdf
% 原页：第 1 页（A3 横排）左栏，印面 921，第 4 题题干下方
% 说明：侧视图为底边 2（左、右各 1）、两腰均为 2 的等腰三角形，
%       高线画虚线；俯视图为左边竖直、右侧为顶点的等腰三角形，
%       底边 2（上、下各 1）、两斜边均为 2，并画出底边中点与
%       右顶点的连线。
% ============================================================
\begin{tikzpicture}[>=Stealth, line width=0.7pt, font=\small,
  x=1cm, y=1cm, line cap=round, line join=round,
  ext/.style={inner sep=2pt}]

  % ---------- 侧视图 ----------
  \draw (0,0) -- (2,0) -- (1,1.732) -- cycle;
  \draw[dashed] (1,0) -- (1,1.732);
  \node[ext,left] at (0.40,0.95) {$2$};
  \node[ext,right] at (1.60,0.95) {$2$};
  \node[ext,below] at (0.5,-0.06) {$1$};
  \node[ext,below] at (1.5,-0.06) {$1$};
  \node at (1,-0.95) {侧视图};

  % ---------- 俯视图 ----------
  \begin{scope}[xshift=4.2cm, yshift=-0.13cm]
    \draw (0,0) -- (0,2) -- (1.732,1) -- cycle;
    \draw (0,1) -- (1.732,1);
    \node[ext,left] at (-0.10,1.55) {$1$};
    \node[ext,left] at (-0.10,0.45) {$1$};
    \node[ext,above left] at (0.90,1.62) {$2$};
    \node[ext,below left] at (0.90,0.38) {$2$};
    \node at (0.87,-0.82) {俯视图};
  \end{scope}
\end{tikzpicture}
\end{center}
\nopagebreak[4]
\Needspace{4\baselineskip}
\textbf{解析}

\Needspace{5\baselineskip}
由俯视图：左边竖直边被标注为 $1$ 与 $1$ 两段，总长 $2$；两条斜边均标注 $2$。故底面是边长为 $2$ 的正三角形，其面积
\nopagebreak[4]
\[S=\frac{\sqrt3}4\times 2^2=\sqrt3 .\]

\Needspace{5\baselineskip}
由侧视图：底边被标注为 $1$ 与 $1$ 两段，总长 $2$，两腰均标注 $2$，故侧视图三角形的高为
\nopagebreak[4]
\[\sqrt{2^2-1^2}=\sqrt3 .\]
侧视图的竖直方向尺寸就是三棱锥的高，故 $h=\sqrt3$。

\Needspace{5\baselineskip}
于是
\nopagebreak[4]
\[V=\frac13Sh=\frac13\times\sqrt3\times\sqrt3=1 .\]

选 D。

\gkanswer{答案：D。俯视图是边长 $2$ 的正三角形，底面积 $S=\sqrt3$；侧视图的高为 $\sqrt3$，即棱锥的高 $h=\sqrt3$，故 $V=\dfrac13Sh=1$。}

\end{answer}
\begin{answer}
\gksolhead{5}{}

\textbf{题目}：

若 $a>b>0$，$c<d<0$，则一定有\mbox{（　　）}

（A）$\dfrac ac>\dfrac bd$；（B）$\dfrac ac<\dfrac bd$；（C）$\dfrac ad>\dfrac bc$；（D）$\dfrac ad<\dfrac bc$。

\Needspace{4\baselineskip}
\textbf{解析}

\Needspace{5\baselineskip}
由 $c<d<0$ 知 $c<0$，$d<0$，从而 $cd>0$。于是
\nopagebreak[4]
\[\frac ad<\frac bc\iff ac<bd\quad(\text{\text{两边同乘} }cd>0).\]
\Needspace{5\baselineskip}
因为 $a>b>0$ 且 $c<0$，所以
\nopagebreak[4]
\[ac<bc ;\]
\Needspace{5\baselineskip}
因为 $b>0$ 且 $c<d$，所以
\nopagebreak[4]
\[bc<bd .\]
两式串联得 $ac<bd$，故 $\dfrac ad<\dfrac bc$ 恒成立。

其余三项均可举反例排除：取 $a=2$，$b=1$，$c=-3$，$d=-1$，则 $\dfrac ac=-\dfrac23$，$\dfrac bd=-1$，$\dfrac ad=-2$，$\dfrac bc=-\dfrac13$，可见（B）$\dfrac ac<\dfrac bd$ 与（C）$\dfrac ad>\dfrac bc$ 都不成立；再取 $a=10$，$b=1$，$c=-1.1$，$d=-1$，则 $\dfrac ac\approx-9.09<\dfrac bd=-1$，可见（A）也不成立。

选 D。

\gkanswer{答案：D。$a/d<b/c\iff ac<bd$；由 $a>b>0$、$c<0$ 得 $ac<bc$，由 $b>0$、$c<d$ 得 $bc<bd$，故 $ac<bd$。}

\end{answer}
\begin{answer}
\gksolhead{6}{}

\textbf{题目}：

执行如图所示的程序框图，如果输入的 $x$，$y\in\mathbf R$，则输出的 $S$ 的最大值为\mbox{（　　）}

（A）$0$；（B）$1$；（C）$2$；（D）$3$。
\nopagebreak[4]
\begin{center}% ============================================================
% f2014_sichuan_arts_q06_flowchart.tex
% 2014 四川卷（文科）第 6 题：程序框图
% 来源 PDF：普通高考/2014/2014四川文.pdf
% 原页：第 1 页（A3 横排）中栏上方，印面 921，第 6 题题干右侧
% 结构：开始 → 输入 x,y → 判断「x≥0，y≥0，x+y≤1？」；
%       是：S=2x+y；否：S=1；两路汇合 → 输出 S → 结束。
% ============================================================
\begin{tikzpicture}[>=Stealth, line width=0.7pt, font=\small,
  x=1cm, y=1cm, line cap=round, line join=round]

  % 开始
  \draw[rounded corners=6pt] (-0.95,0.45) rectangle (0.95,-0.45);
  \node at (0,0) {开始};
  \draw[->] (0,-0.45) -- (0,-0.90);

  % 输入
  \draw (-1.20,-0.90) -- (1.70,-0.90) -- (1.20,-1.60) -- (-1.70,-1.60) -- cycle;
  \node at (0,-1.25) {输入 $x$，$y$};
  \draw[->] (0,-1.60) -- (0,-2.00);

  % 判断（六边形，与原页一致）
  \draw (-2.90,-2.90) -- (-2.30,-2.40) -- (2.30,-2.40) -- (2.90,-2.90)
     -- (2.30,-3.40) -- (-2.30,-3.40) -- cycle;
  \node at (0,-2.78) {$x\geqslant 0$，$y\geqslant 0$，};
  \node at (0,-3.12) {$x+y\leqslant 1$？};

  % 是：左支
  \draw[->] (-2.90,-2.90) -- (-3.90,-2.90) -- (-3.90,-3.70);
  \node at (-3.15,-2.62) {是};
  \draw (-5.20,-3.70) rectangle (-2.60,-4.50);
  \node at (-3.90,-4.10) {$S=2x+y$};

  % 否：右支
  \draw[->] (2.90,-2.90) -- (3.90,-2.90) -- (3.90,-3.70);
  \node at (3.15,-2.62) {否};
  \draw (2.60,-3.70) rectangle (5.20,-4.50);
  \node at (3.90,-4.10) {$S=1$};

  % 两路汇合
  \draw (-3.90,-4.50) -- (-3.90,-4.90) -- (3.90,-4.90) -- (3.90,-4.50);
  \draw[->] (0,-4.90) -- (0,-5.25);

  % 输出
  \draw (-1.10,-5.25) -- (1.80,-5.25) -- (1.30,-5.95) -- (-1.60,-5.95) -- cycle;
  \node at (0.10,-5.60) {输出 $S$};
  \draw[->] (0.10,-5.95) -- (0.10,-6.40);

  % 结束
  \draw[rounded corners=6pt] (-0.85,-6.40) rectangle (1.05,-7.30);
  \node at (0.10,-6.85) {结束};
\end{tikzpicture}
\end{center}
\nopagebreak[4]
\Needspace{4\baselineskip}
\textbf{解析}

\Needspace{5\baselineskip}
当 $x\geqslant 0$，$y\geqslant 0$，$x+y\leqslant 1$ 时，$S=2x+y$。由 $x+y\leqslant 1$ 得
\nopagebreak[4]
\[2x+y=x+(x+y)\leqslant x+1\leqslant 2,\]
且当 $x=1$，$y=0$ 时 $S=2$，故此时 $S$ 的最大值为 $2$。

\Needspace{5\baselineskip}
当条件不成立时 $S=1$。比较两种情形，
\nopagebreak[4]
\[\max S=\max\{2,1\}=2 .\]

选 C。

\gkanswer{答案：C。条件成立时 $S=2x+y$ 在 $x=1$，$y=0$ 处取最大值 $2$，否则 $S=1$，故 $S$ 的最大值为 $2$。}

\end{answer}
\begin{answer}
\gksolhead{7}{}

\textbf{题目}：

已知 $b>0$，$\log_5 b=a$，$\lg b=c$，$5^d=10$，则下列等式一定成立的是\mbox{（　　）}

（A）$d=ac$；（B）$a=cd$；（C）$c=ad$；（D）$d=a+c$。

\Needspace{4\baselineskip}
\textbf{解析}

\Needspace{5\baselineskip}
由换底公式，
\nopagebreak[4]
\[a=\log_5 b=\frac{\lg b}{\lg 5}=\frac{c}{\lg 5}.\]
\Needspace{5\baselineskip}
又 $5^d=10$ 给出
\nopagebreak[4]
\[d=\log_5 10=\frac{1}{\lg 5}.\]
\Needspace{5\baselineskip}
于是
\nopagebreak[4]
\[cd=\lg b\cdot\frac1{\lg 5}=\frac{c}{\lg 5}=a .\]

逐项检验，取 $b=5$，则 $a=1$，$c=\lg5\approx0.699$，$d=\log_510\approx1.431$：（A）$ac\approx0.699\neq d$；（C）$ad\approx1.431\neq c$；（D）$a+c\approx1.699\neq d$。只有（B）对任意 $b>0$ 恒成立。

选 B。

\gkanswer{答案：B。$a=\log_5b=\dfrac{\lg b}{\lg 5}$，$d=\log_510=\dfrac1{\lg 5}$，故 $cd=\dfrac{\lg b}{\lg5}=a$。}

\end{answer}
\begin{answer}
\gksolhead{8}{}

\textbf{题目}：

如图，从气球 $A$ 上测得正前方的河流的两岸 $B$，$C$ 的俯角分别为 $75^\circ$，$30^\circ$，此时气球的高是 $60\ \mathrm m$，则河流的宽度 $BC$ 等于\mbox{（　　）}

（A）$240(\sqrt3-1)\ \mathrm m$；（B）$180(\sqrt2-1)\ \mathrm m$；（C）$120(\sqrt3-1)\ \mathrm m$；（D）$30(\sqrt3+1)\ \mathrm m$。
\nopagebreak[4]
\begin{center}% ============================================================
% f2014_sichuan_arts_q08_river_width.tex
% 2014 四川卷（文科）第 8 题：气球上测河流宽度
% 来源 PDF：普通高考/2014/2014四川文.pdf
% 原页：第 1 页（A3 横排）中栏，印面 921，第 8 题题干下方
% 说明：A 为气球位置，其正下方为河岸上的垂足；B、C 为对岸与
%       本岸的两个观测点，A 到 B、C 的俯角分别为 75°、30°；
%       气球高 60 m，图中给出水平虚线、铅垂高度线与直角标记。
% ============================================================
\begin{tikzpicture}[>=Stealth, line width=0.7pt, font=\small,
  x=1cm, y=1cm, line cap=round, line join=round]

  \coordinate (A) at (0,4);
  \coordinate (D) at (0,0);
  \coordinate (B) at (1.10,0);
  \coordinate (C) at (6.90,0);

  % 地面与过 A 的水平线
  \draw[dashed] (-0.80,0) -- (7.60,0);
  \draw[dashed] (-0.20,4) -- (2.20,4);

  % 视线
  \draw[dashed] (A) -- (B);
  \draw[dashed] (A) -- (C);

  % 气球
  \draw (0,4.44) circle (0.20);
  \draw (-0.08,4.26) -- (0.08,4.26) -- (0.10,4.00) -- (-0.10,4.00) -- cycle;
  \node at (0,4.45) {$A$};

  % 铅垂高度与标注
  \draw (-0.85,4) -- (-0.25,4);
  \draw (-0.85,0) -- (-0.25,0);
  \draw[<->] (-0.55,4) -- (-0.55,0);
  \node[left] at (-0.70,2) {$60\mathrm{m}$};

  % 垂足处的直角标记
  \draw (0,0.16) -- (0.16,0.16) -- (0.16,0);

  % 俯角
  \draw ($(A)+(0.60,0)$) arc (0:-30:0.60);
  \node at (0.98,3.80) {$30^\circ$};
  \draw ($(A)+(1.20,0)$) arc (0:-75:1.20);
  \node at (1.78,3.62) {$75^\circ$};

  % 点的标注
  \node[below] at (B) {$B$};
  \node[below] at (C) {$C$};
\end{tikzpicture}
\end{center}
\nopagebreak[4]
\Needspace{4\baselineskip}
\textbf{解析}

\Needspace{5\baselineskip}
设气球 $A$ 正下方的河岸点为 $D$，则 $AD=60\ \mathrm m$，且 $AD\perp BD$，$AD\perp CD$。过 $A$ 的水平线与河岸平行，由俯角的定义及内错角相等，
\nopagebreak[4]
\[\angle ABD=75^\circ,\qquad \angle ACD=30^\circ .\]

\Needspace{5\baselineskip}
在 $\mathrm{Rt}\triangle ABD$ 中，
\nopagebreak[4]
\[BD=\frac{AD}{\tan75^\circ}=60(2-\sqrt3)\ \mathrm m ;\]
\Needspace{5\baselineskip}
在 $\mathrm{Rt}\triangle ACD$ 中，
\nopagebreak[4]
\[CD=\frac{AD}{\tan30^\circ}=60\sqrt3\ \mathrm m .\]
\Needspace{5\baselineskip}
$B$ 在 $D$，$C$ 之间，故
\nopagebreak[4]
\[BC=CD-BD=60\sqrt3-60(2-\sqrt3)=120(\sqrt3-1)\ \mathrm m .\]

选 C。

\gkanswer{答案：C。设垂足为 $D$，$AD=60$，则 $BD=60\cot75^\circ=60(2-\sqrt3)$，$CD=60\cot30^\circ=60\sqrt3$，故 $BC=CD-BD=120(\sqrt3-1)\ \mathrm m$。}

\end{answer}
\begin{answer}
\gksolhead{9}{}

\textbf{题目}：

设 $m\in\mathbf R$，过定点 $A$ 的动直线 $x+my=0$ 和过定点 $B$ 的动直线 $mx-y-m+3=0$ 交于点 $P(x,y)$，则 $|PA|+|PB|$ 的取值范围是\mbox{（　　）}

（A）$[\sqrt5,2\sqrt5]$；（B）$[\sqrt{10},2\sqrt5]$；（C）$[\sqrt{10},4\sqrt5]$；（D）$[2\sqrt5,4\sqrt5]$。

\Needspace{4\baselineskip}
\textbf{解析}

直线 $x+my=0$ 恒过原点，故 $A(0,0)$。

\Needspace{5\baselineskip}
直线 $mx-y-m+3=0$ 即
\nopagebreak[4]
\[m(x-1)-(y-3)=0,\]
恒过点 $(1,3)$，故 $B(1,3)$。

\Needspace{5\baselineskip}
两直线的斜率分别为 $-\dfrac1m$（$m\neq0$）与 $m$，乘积为 $-1$；当 $m=0$ 时两直线为 $x=0$ 与 $y=3$，也互相垂直。总之恒有 $PA\perp PB$，故点 $P$ 的轨迹是以 $AB$ 为直径的圆（$m=3$ 时 $P=A$，$m=-\dfrac13$ 时 $P=B$，是两个端点情形），由勾股定理
\nopagebreak[4]
\[|PA|^2+|PB|^2=|AB|^2=1^2+3^2=10 .\]
\Needspace{5\baselineskip}
于是
\nopagebreak[4]
\[\bigl(|PA|+|PB|\bigr)^2=10+2|PA|\cdot|PB|,\]
\Needspace{5\baselineskip}
而由 $2|PA|\cdot|PB|\leqslant |PA|^2+|PB|^2=10$（等号当 $|PA|=|PB|=\sqrt5$）得
\nopagebreak[4]
\[10\leqslant\bigl(|PA|+|PB|\bigr)^2\leqslant 20,\]
\Needspace{5\baselineskip}
即
\nopagebreak[4]
\[\sqrt{10}\leqslant|PA|+|PB|\leqslant 2\sqrt5 .\]

端点可达到：$m=3$ 时两直线为 $x+3y=0$ 与 $3x-y=0$，交于 $P=A$，此时 $|PA|+|PB|=\sqrt{10}$；$m=-\dfrac13$ 时交于 $P=B$；$|PA|=|PB|=\sqrt5$ 时取到 $2\sqrt5$。

选 B。

\gkanswer{答案：B。$A(0,0)$，$B(1,3)$，两直线垂直，故 $P$ 在以 $AB$ 为直径的圆上，$|PA|^2+|PB|^2=10$，于是 $|PA|+|PB|\in[\sqrt{10},2\sqrt5]$。}

\end{answer}
\begin{answer}
\gksolhead{10}{}

\textbf{题目}：

已知 $F$ 是抛物线 $y^2=x$ 的焦点，点 $A$，$B$ 在该抛物线上且位于 $x$ 轴的两侧，$\overrightarrow{OA}\cdot\overrightarrow{OB}=2$（其中 $O$ 为坐标原点），则 $\triangle ABO$ 与 $\triangle AFO$ 面积之和的最小值是\mbox{（　　）}

（A）$2$；（B）$3$；（C）$\dfrac{17\sqrt2}8$；（D）$\sqrt{10}$。

\Needspace{4\baselineskip}
\textbf{解析}

\Needspace{5\baselineskip}
抛物线 $y^2=x$ 的焦点为 $F\left(\dfrac14,0\right)$。设
\nopagebreak[4]
\[A(y_1^2,y_1),\qquad B(y_2^2,y_2),\]
$A$，$B$ 位于 $x$ 轴两侧，不妨设 $y_1>0>y_2$，则 $y_1y_2<0$。

\Needspace{5\baselineskip}
由 $\overrightarrow{OA}\cdot\overrightarrow{OB}=2$，
\nopagebreak[4]
\[y_1^2y_2^2+y_1y_2=2 .\]
\Needspace{5\baselineskip}
令 $u=y_1y_2$，得 $u^2+u-2=0$，即 $u=1$ 或 $u=-2$；结合 $u<0$ 得
\nopagebreak[4]
\[y_1y_2=-2,\qquad y_2=-\frac2{y_1}\ (y_1>0).\]

\Needspace{5\baselineskip}
三角形面积
\nopagebreak[4]
\[S_{\triangle ABO}=\frac12\left|y_1^2y_2-y_2^2y_1\right|
=\frac12|y_1y_2|\cdot|y_1-y_2|=|y_1-y_2|=y_1-y_2 .\]
\Needspace{5\baselineskip}
又 $|OF|=\dfrac14$，故
\nopagebreak[4]
\[S_{\triangle AFO}=\frac12\cdot\frac14\cdot y_1=\frac{y_1}8 .\]
\Needspace{5\baselineskip}
面积之和
\nopagebreak[4]
\[S=y_1-y_2+\frac{y_1}8=\frac98y_1+\frac2{y_1}\geqslant 2\sqrt{\frac98\times2}=3,\]
等号当 $\dfrac98y_1=\dfrac2{y_1}$，即 $y_1^2=\dfrac{16}9$、$y_1=\dfrac43$ 时成立。

故最小值为 $3$。

选 B。

\gkanswer{答案：B。设 $A(y_1^2,y_1)$，$B(y_2^2,y_2)$，由 $\overrightarrow{OA}\cdot\overrightarrow{OB}=2$ 及 $y_1y_2<0$ 得 $y_1y_2=-2$；面积之和为 $\dfrac98y_1-y_2=\dfrac98y_1+\dfrac2{y_1}\geqslant 3$，当 $y_1=\dfrac43$ 时取等号。}

\end{answer}
\begin{answer}
\gksolhead{11}{}

\textbf{题目}：

双曲线 $\dfrac{x^2}4-y^2=1$ 的离心率等于\gkblank{}。

\Needspace{4\baselineskip}
\textbf{解析}

\Needspace{5\baselineskip}
双曲线 $\dfrac{x^2}4-y^2=1$ 中 $a^2=4$，$b^2=1$，故
\nopagebreak[4]
\[c=\sqrt{a^2+b^2}=\sqrt5,\qquad e=\frac ca=\frac{\sqrt5}2 .\]

\gkanswer{答案：$\dfrac{\sqrt5}2$。$a=2$，$b=1$，$c=\sqrt{a^2+b^2}=\sqrt5$，故 $e=\dfrac ca=\dfrac{\sqrt5}2$。}

\end{answer}
\begin{answer}
\gksolhead{12}{}

\textbf{题目}：

复数 $\dfrac{2-2\mathrm i}{1+\mathrm i}=$\gkblank{}。

\Needspace{4\baselineskip}
\textbf{解析}

\Needspace{5\baselineskip}
分子分母同乘 $1-\mathrm i$：
\nopagebreak[4]
\[\frac{2-2\mathrm i}{1+\mathrm i}
=\frac{(2-2\mathrm i)(1-\mathrm i)}{(1+\mathrm i)(1-\mathrm i)}
=\frac{2-2\mathrm i-2\mathrm i+2\mathrm i^2}{2}
=\frac{2-4\mathrm i-2}{2}=-2\mathrm i .\]

\gkanswer{答案：$-2\mathrm i$。分子分母同乘 $1-\mathrm i$，得 $\dfrac{(2-2\mathrm i)(1-\mathrm i)}2=\dfrac{-4\mathrm i}2=-2\mathrm i$。}

\end{answer}
\begin{answer}
\gksolhead{13}{}

\textbf{题目}：

\Needspace{5\baselineskip}
设 $f(x)$ 是定义在 $\mathbf R$ 上的周期为 $2$ 的函数，当 $x\in[-1,1)$ 时，
\nopagebreak[4]
\[f(x)=\begin{cases}-4x^2+2,&-1\leqslant x<0,\\x,&0\leqslant x<1,\end{cases}\]
则 $f\left(\dfrac32\right)=$\gkblank{}。

\Needspace{4\baselineskip}
\textbf{解析}

\Needspace{5\baselineskip}
$f(x)$ 的周期为 $2$，故
\nopagebreak[4]
\[f\left(\frac32\right)=f\left(\frac32-2\right)=f\left(-\frac12\right).\]
\Needspace{5\baselineskip}
而 $-\dfrac12\in[-1,0)$，用第一段解析式：
\nopagebreak[4]
\[f\left(-\frac12\right)=-4\times\left(\frac12\right)^2+2=-1+2=1 .\]

\gkanswer{答案：$1$。周期为 $2$，故 $f\left(\dfrac32\right)=f\left(-\dfrac12\right)=-4\times\dfrac14+2=1$。}

\end{answer}
\begin{answer}
\gksolhead{14}{}

\textbf{题目}：

平面向量 $\boldsymbol a=(1,2)$，$\boldsymbol b=(4,2)$，$\boldsymbol c=m\boldsymbol a+\boldsymbol b$（$m\in\mathbf R$），且 $\boldsymbol c$ 与 $\boldsymbol a$ 的夹角等于 $\boldsymbol c$ 与 $\boldsymbol b$ 的夹角，则 $m=$\gkblank{}。

\Needspace{4\baselineskip}
\textbf{解析}

\Needspace{5\baselineskip}
先算
\nopagebreak[4]
\[\boldsymbol c=m(1,2)+(4,2)=(m+4,\,2m+2).\]
若 $\boldsymbol c=\boldsymbol 0$，则 $m+4=0$ 且 $2m+2=0$，无解，故 $\boldsymbol c\neq\boldsymbol 0$。

\Needspace{5\baselineskip}
夹角相等即余弦相等，而 $|\boldsymbol c|$ 在两式中相同，故
\nopagebreak[4]
\[\frac{\boldsymbol c\cdot\boldsymbol a}{|\boldsymbol a|}=\frac{\boldsymbol c\cdot\boldsymbol b}{|\boldsymbol b|} .\]
\Needspace{5\baselineskip}
逐项计算：
\nopagebreak[4]
\[\boldsymbol c\cdot\boldsymbol a=(m+4)+2(2m+2)=5m+8,\qquad |\boldsymbol a|=\sqrt5,\]
\nopagebreak[4]
\[\boldsymbol c\cdot\boldsymbol b=4(m+4)+2(2m+2)=8m+20,\qquad |\boldsymbol b|=\sqrt{20}=2\sqrt5 .\]
\Needspace{5\baselineskip}
代入得
\nopagebreak[4]
\[\frac{5m+8}{\sqrt5}=\frac{8m+20}{2\sqrt5}
\ \Longrightarrow\ 2(5m+8)=8m+20
\ \Longrightarrow\ m=2 .\]

\gkanswer{答案：$2$。$\boldsymbol c=(m+4,2m+2)$，由两夹角相等得 $\dfrac{5m+8}{\sqrt5}=\dfrac{8m+20}{2\sqrt5}$，解得 $m=2$。}

\end{answer}
\begin{answer}
\gksolhead{15}{}

\textbf{题目}：

以 $A$ 表示值域为 $\mathbf R$ 的函数组成的集合，$B$ 表示具有如下性质的函数 $\varphi(x)$ 组成的集合：对于函数 $\varphi(x)$，存在一个正数 $M$，使得函数 $\varphi(x)$ 的值域包含于区间 $[-M,M]$。例如，当 $\varphi_1(x)=x^3$，$\varphi_2(x)=\sin x$ 时，$\varphi_1(x)\in A$，$\varphi_2(x)\in B$。现有如下命题：

① 设函数 $f(x)$ 的定义域为 $D$，则“$f(x)\in A$”的充要条件是“$\forall b\in\mathbf R$，$\exists a\in D$，$f(a)=b$”；

② 函数 $f(x)\in B$ 的充要条件是 $f(x)$ 有最大值和最小值；

③ 若函数 $f(x)$，$g(x)$ 的定义域相同，且 $f(x)\in A$，$g(x)\in B$，则 $f(x)+g(x)\notin B$；

④ 若函数 $f(x)=a\ln(x+2)+\dfrac{x}{x^2+1}$（$x>-2$，$a\in\mathbf R$）有最大值，则 $f(x)\in B$。

其中的真命题有\gkblank{}。（写出所有真命题的序号）

\Needspace{4\baselineskip}
\textbf{解析}

$A$ 是值域为 $\mathbf R$ 的函数之集，$B$ 是值域有界的函数之集，且二者互不相交。

① 值域为 $\mathbf R$ 的定义就是：对任意 $b\in\mathbf R$，存在 $a\in D$ 使 $f(a)=b$。故①真。

② 有最大值和最小值只能推出值域有界，反之不真：$f(x)=\arctan x$ 的值域为 $\left(-\dfrac\pi2,\dfrac\pi2\right)\subset[-2,2]$，故 $f\in B$，但它既无最大值也无最小值。故②假。

\Needspace{5\baselineskip}
③ 若 $f(x)+g(x)\in B$，则存在 $M_1>0$ 使 $|f(x)+g(x)|\leqslant M_1$；又 $g\in B$，存在 $M_2>0$ 使 $|g(x)|\leqslant M_2$。于是
\nopagebreak[4]
\[|f(x)|\leqslant|f(x)+g(x)|+|g(x)|\leqslant M_1+M_2,\]
即 $f$ 的值域有界，与 $f\in A$ 矛盾。故③真。

\Needspace{5\baselineskip}
④ 若 $a>0$，则 $x\to+\infty$ 时 $a\ln(x+2)\to+\infty$，$f(x)\to+\infty$，无最大值；若 $a<0$，则 $x\to-2^+$ 时 $a\ln(x+2)\to+\infty$，同样无最大值。故只能 $a=0$，此时
\nopagebreak[4]
\[f(x)=\frac{x}{x^2+1},\qquad |f(x)|=\frac{|x|}{x^2+1}\leqslant\frac12 ,\]
等号在 $x=\pm1$ 处成立，值域为 $\left[-\dfrac12,\dfrac12\right]$，故 $f\in B$。④真。

综上，真命题为①③④。

\gkanswer{答案：①③④。①为满射的等价表述；②错，如 $\arctan x$ 有界但无最值；③对，否则 $f=(f+g)-g$ 也有界；④对，$a>0$ 或 $a<0$ 时 $f$ 无最大值，只能 $a=0$，此时 $f=\dfrac{x}{x^2+1}$ 有最大值和最小值。}

\end{answer}
\begin{answer}
\gksolhead{16}{}

\textbf{题目}：

一个盒子里装有三张卡片，分别标记有数字 $1,2,3$，这三张卡片除标记的数字外完全相同。随机有放回地抽取 $3$ 次，每次抽取 $1$ 张，将抽取的卡片上的数字依次记为 $a$，$b$，$c$。

（1） 求“抽取的卡片上的数字满足 $a+b=c$”的概率；

（2） 求“抽取的卡片上的数字 $a$，$b$，$c$ 不完全相同”的概率。

\Needspace{4\baselineskip}
\textbf{解析}

\Needspace{5\baselineskip}
有放回地抽取 $3$ 次，每次有 $3$ 种可能，故基本事件总数为
\nopagebreak[4]
\[3^3=27,\]
且每个有序三元组 $(a,b,c)$ 等可能。

\gksubpart{（1）}

\Needspace{5\baselineskip}
逐一列举满足 $a+b=c$ 的三元组：$a,b,c\in\{1,2,3\}$ 且 $c\leqslant 3$，故 $a+b\leqslant 3$。只有
\nopagebreak[4]
\[(1,1,2),\quad(1,2,3),\quad(2,1,3)\]
\Needspace{5\baselineskip}
共 $3$ 个。所求概率为
\nopagebreak[4]
\[P=\frac3{27}=\frac19 .\]

\gksubpart{（2）}

\Needspace{5\baselineskip}
“$a$，$b$，$c$ 不完全相同”的对立事件是“$a=b=c$”，后者只有
\nopagebreak[4]
\[(1,1,1),\quad(2,2,2),\quad(3,3,3)\]
\Needspace{5\baselineskip}
共 $3$ 个。故所求概率为
\nopagebreak[4]
\[P=1-\frac3{27}=\frac{24}{27}=\frac89 .\]

\gkanswer{答案（1）：$\dfrac19$。有放回抽取 $3$ 次共 $27$ 个等可能结果，满足 $a+b=c$ 的只有 $(1,1,2)$，$(1,2,3)$，$(2,1,3)$ 三个。}

\gkanswer{答案（2）：$\dfrac89$。$a$，$b$，$c$ 完全相同的只有 $(1,1,1)$，$(2,2,2)$，$(3,3,3)$ 三个，故所求为 $1-\dfrac3{27}=\dfrac89$。}

\end{answer}
\begin{answer}
\gksolhead{17}{}

\textbf{题目}：

已知函数 $f(x)=\sin\left(3x+\dfrac\pi4\right)$。

（1） 求 $f(x)$ 的单调递增区间；

（2） 若 $\alpha$ 是第二象限角，$f\left(\dfrac\alpha3\right)=\dfrac45\cos\left(\alpha+\dfrac\pi4\right)\cos 2\alpha$，求 $\cos\alpha-\sin\alpha$ 的值。

\Needspace{4\baselineskip}
\textbf{解析}

\gksubpart{（1）}

\Needspace{5\baselineskip}
正弦函数 $\sin t$ 在
\nopagebreak[4]
\[-\frac\pi2+2k\pi\leqslant t\leqslant\frac\pi2+2k\pi\quad(k\in\mathbf Z)\]
\Needspace{5\baselineskip}
上单调递增。令 $t=3x+\dfrac\pi4$，得
\nopagebreak[4]
\[-\frac\pi2+2k\pi\leqslant 3x+\frac\pi4\leqslant\frac\pi2+2k\pi,\]
\Needspace{5\baselineskip}
即
\nopagebreak[4]
\[-\frac{3\pi}4+2k\pi\leqslant 3x\leqslant\frac\pi4+2k\pi,\]
\nopagebreak[4]
\[-\frac\pi4+\frac{2k\pi}3\leqslant x\leqslant\frac\pi{12}+\frac{2k\pi}3 .\]
\Needspace{5\baselineskip}
故 $f(x)$ 的单调递增区间为
\nopagebreak[4]
\[\left[-\frac\pi4+\frac{2k\pi}3,\ \frac\pi{12}+\frac{2k\pi}3\right]\quad(k\in\mathbf Z).\]

\gksubpart{（2）}

\Needspace{5\baselineskip}
由 $f\left(\dfrac\alpha3\right)=\sin\left(\alpha+\dfrac\pi4\right)$，题设条件即
\nopagebreak[4]
\[\sin\left(\alpha+\frac\pi4\right)=\frac45\cos\left(\alpha+\frac\pi4\right)\cos2\alpha .\]
\Needspace{5\baselineskip}
记 $s=\sin\alpha$，$c=\cos\alpha$，则
\nopagebreak[4]
\[\sin\left(\alpha+\frac\pi4\right)=\frac{\sqrt2}2(s+c),\qquad
\cos\left(\alpha+\frac\pi4\right)=\frac{\sqrt2}2(c-s),\]
\nopagebreak[4]
\[\cos2\alpha=c^2-s^2=(c-s)(c+s).\]
\Needspace{5\baselineskip}
代入得
\nopagebreak[4]
\[\frac{\sqrt2}2(s+c)=\frac45\cdot\frac{\sqrt2}2(c-s)^2(s+c).\]

\Needspace{5\baselineskip}
若 $s+c=0$，由 $\alpha$ 在第二象限（$s>0,c<0$）及 $s^2+c^2=1$ 得
\nopagebreak[4]
\[\sin\alpha=\frac{\sqrt2}2,\qquad \cos\alpha=-\frac{\sqrt2}2,\qquad \alpha=\frac{3\pi}4,\]
此时 $\cos\alpha-\sin\alpha=-\sqrt2$。

\Needspace{5\baselineskip}
若 $s+c\neq0$，约去 $\frac{\sqrt2}2(s+c)$ 得
\nopagebreak[4]
\[(c-s)^2=\frac54 .\]
\Needspace{5\baselineskip}
由 $|c-s|\leqslant\sqrt2$ 知 $|c-s|=\frac{\sqrt5}2<\sqrt2$ 可行；又 $\alpha$ 在第二象限，$c-s<0$，故 $c-s=-\frac{\sqrt5}2$，即
\nopagebreak[4]
\[\cos\alpha-\sin\alpha=-\frac{\sqrt5}2 .\]
此时 $\cos\left(\alpha+\frac\pi4\right)=\frac{c-s}{\sqrt2}=-\frac{\sqrt{10}}4$，$\alpha\approx97.2^\circ$ 或 $172.8^\circ$，均属第二象限。

所以原题条件下 $\cos\alpha-\sin\alpha$ 有两个可能值：$-\sqrt2$ 或 $-\dfrac{\sqrt5}2$；原题未给出排除任一支的条件。

\gkanswer{答案（1）：$\left[-\dfrac\pi4+\dfrac{2k\pi}3,\ \dfrac\pi{12}+\dfrac{2k\pi}3\right]$（$k\in\mathbf Z$）。由 $-\dfrac\pi2+2k\pi\leqslant 3x+\dfrac\pi4\leqslant\dfrac\pi2+2k\pi$ 解出。}

\gkanswer{答案（2）：$-\sqrt2$ 或 $-\dfrac{\sqrt5}2$。记 $u=\cos\alpha-\sin\alpha$、$v=\cos\alpha+\sin\alpha$，条件化为 $v=\dfrac45u^2v$：$v=0$ 给 $\alpha=\dfrac{3\pi}4$、$u=-\sqrt2$；$u^2=\dfrac54$ 给 $u=-\dfrac{\sqrt5}2$（$\alpha\approx97.2^\circ$ 或 $172.8^\circ$），两支均满足第二象限条件。}

\end{answer}
\begin{answer}
\gksolhead{18}{}

\textbf{题目}：

在如图所示的多面体中，四边形 $ABB_1A_1$ 和 $ACC_1A_1$ 都为矩形。
\nopagebreak[4]
\begin{center}% ============================================================
% f2014_sichuan_arts_q18_prism.tex
% 2014 四川卷（文科）第 18 题：多面体（三棱柱 ABC-A1B1C1）
% 来源 PDF：普通高考/2014/2014四川文.pdf
% 原页：第 2 页（A3 横排）左栏，印面 922，第 18 题题干右侧
% 说明：四边形 ABB1A1 与 ACC1A1 均为矩形，故 AA1、BB1、CC1
%       两两平行且都垂直于底面；AC 为被遮挡的棱，画虚线。
%       D 为 BC 中点，E 为 CC1 中点，DE 为辅助线。
% ============================================================
\begin{tikzpicture}[>=Stealth, line width=0.7pt, font=\small,
  x=1cm, y=1cm, line cap=round, line join=round]

  \coordinate (A)  at (0,0);
  \coordinate (B)  at (2.74,-1.65);
  \coordinate (C)  at (4.30,0);
  \coordinate (A1) at (0,3.24);
  \coordinate (B1) at (2.74,1.59);
  \coordinate (C1) at (4.30,3.24);
  \coordinate (D)  at (3.52,-0.825);
  \coordinate (E)  at (4.30,1.62);

  % 被遮挡的棱
  \draw[dashed] (A) -- (C);

  % 可见的棱
  \draw (A) -- (B) -- (C) -- (C1) -- (A1) -- cycle;
  \draw (A1) -- (B1) -- (C1);
  \draw (B1) -- (B);

  % 辅助线 DE
  \draw (D) -- (E);

  % 顶点标注
  \node[left] at (A1) {$A_1$};
  \node[above right] at (C1) {$C_1$};
  \node[left] at (A) {$A$};
  \node[right] at (C) {$C$};
  \node[below right] at (B1) {$B_1$};
  \node[right] at (E) {$E$};
  \node[below right] at (D) {$D$};
  \node[below] at (B) {$B$};
\end{tikzpicture}
\end{center}
\nopagebreak[4]
（1） 若 $AC\perp BC$，证明：直线 $BC\perp$ 平面 $ACC_1A_1$；

（2） 设 $D$，$E$ 分别是线段 $BC$，$CC_1$ 的中点，在线段 $AB$ 上是否存在一点 $M$，使直线 $DE\parallel$ 平面 $A_1MC$？请证明你的结论。

\Needspace{4\baselineskip}
\textbf{解析}

\gksubpart{（1）}

\Needspace{5\baselineskip}
因为四边形 $ACC_1A_1$ 是矩形，所以
\nopagebreak[4]
\[CC_1\perp AC .\]
\Needspace{5\baselineskip}
又已知 $BC\perp AC$，且 $AC\cap CC_1=C$，$AC$，$CC_1$ 都在平面 $ACC_1A_1$ 内，由线面垂直的判定定理，
\nopagebreak[4]
\[BC\perp\text{\text{平面} }ACC_1A_1 .\]

\gksubpart{（2）}

\textbf{结论：存在，且 $M$ 为 $AB$ 的中点。}

\Needspace{5\baselineskip}
因为 $ABB_1A_1$ 与 $ACC_1A_1$ 都是矩形，所以
\nopagebreak[4]
\[AA_1\perp AB,\qquad AA_1\perp AC .\]
于是 $AA_1\perp$ 平面 $ABC$，而 $AB$ 与 $AC$ 不共线，故 $\overrightarrow{AB}$，$\overrightarrow{AC}$，$\overrightarrow{AA_1}$ 不共面，可作为空间的一组基。
\Needspace{5\baselineskip}
令 $\overrightarrow{AB}=\boldsymbol u$，$\overrightarrow{AC}=\boldsymbol v$，$\overrightarrow{AA_1}=\boldsymbol w$，以 $A$ 为原点，则
\nopagebreak[4]
\[B=\boldsymbol u,\quad C=\boldsymbol v,\quad A_1=\boldsymbol w,\quad B_1=\boldsymbol u+\boldsymbol w .\]
\Needspace{5\baselineskip}
$D$，$E$ 分别是 $BC$，$CC_1$ 的中点，故
\nopagebreak[4]
\[D=\frac{\boldsymbol u+\boldsymbol v}2,\qquad E=\boldsymbol v+\frac{\boldsymbol w}2,\]
\nopagebreak[4]
\[\overrightarrow{DE}=E-D=-\frac{\boldsymbol u}2+\frac{\boldsymbol v}2+\frac{\boldsymbol w}2 .\]

\Needspace{5\baselineskip}
设 $M=t\boldsymbol u$（$0\leqslant t\leqslant 1$），则平面 $A_1MC$ 由向量
\nopagebreak[4]
\[\overrightarrow{A_1M}=t\boldsymbol u-\boldsymbol w,\qquad \overrightarrow{A_1C}=\boldsymbol v-\boldsymbol w\]
\Needspace{5\baselineskip}
张成。若 $DE\parallel$ 平面 $A_1MC$，则存在实数 $\alpha$，$\beta$ 使
\nopagebreak[4]
\[\overrightarrow{DE}=\alpha(t\boldsymbol u-\boldsymbol w)+\beta(\boldsymbol v-\boldsymbol w).\]
\Needspace{5\baselineskip}
比较 $\boldsymbol u$，$\boldsymbol v$，$\boldsymbol w$ 的系数：
\nopagebreak[4]
\[\alpha t=-\frac12,\qquad \beta=\frac12,\qquad -(\alpha+\beta)=\frac12 .\]
\Needspace{5\baselineskip}
由第三式得 $\alpha=-1$，代入第一式得
\nopagebreak[4]
\[t=\frac12 .\]
所以 $M$ 只能是 $AB$ 的中点。

\Needspace{5\baselineskip}
反之，取 $M$ 为 $AB$ 的中点（$t=\dfrac12$，$\alpha=-1$，$\beta=\dfrac12$），有
\nopagebreak[4]
\[\overrightarrow{DE}=-\overrightarrow{A_1M}+\frac12\overrightarrow{A_1C},\]
\Needspace{5\baselineskip}
即 $\overrightarrow{DE}$ 是平面 $A_1MC$ 内两个向量的线性组合，故 $DE\parallel$ 平面 $A_1MC$。又若 $E$ 在平面 $A_1MC$ 内，则存在 $s$，$r$ 使
\nopagebreak[4]
\[\boldsymbol v+\frac{\boldsymbol w}2=\boldsymbol w+s\left(\frac{\boldsymbol u}2-\boldsymbol w\right)+r(\boldsymbol v-\boldsymbol w),\]
比较系数得 $s=0$，$r=1$，$\dfrac12=0$，矛盾。故 $E$ 不在平面内，$DE$ 不在平面内。

所以存在满足条件的点 $M$，它就是 $AB$ 的中点。

\gkanswer{答案（1）：证明。矩形 $ACC_1A_1$ 给出 $CC_1\perp AC$，又 $BC\perp AC$，且 $AC\cap CC_1=C$，故 $BC\perp$ 平面 $ACC_1A_1$。}

\gkanswer{答案（2）：存在，$M$ 为 $AB$ 的中点。以 $\overrightarrow{AB}$，$\overrightarrow{AC}$，$\overrightarrow{AA_1}$ 为基表示，$M$ 为中点时 $\overrightarrow{DE}=-\overrightarrow{A_1M}+\dfrac12\overrightarrow{A_1C}$，故 $DE\parallel$ 平面 $A_1MC$；由系数比较知这样的 $M$ 唯一。}

\end{answer}
\begin{answer}
\gksolhead{19}{}

\textbf{题目}：

设等差数列 $\{a_n\}$ 的公差为 $d$，点 $(a_n,b_n)$ 在函数 $f(x)=2^x$ 的图象上（$n\in\mathbf N^*$）。

（1） 证明：数列 $\{b_n\}$ 为等比数列；

（2） 若 $a_1=1$，函数 $f(x)$ 的图象在点 $(a_2,b_2)$ 处的切线在 $x$ 轴上的截距为 $2-\dfrac1{\ln 2}$，求数列 $\{a_nb_n^2\}$ 的前 $n$ 项和 $S_n$。

\Needspace{4\baselineskip}
\textbf{解析}

\gksubpart{（1）}

\Needspace{5\baselineskip}
由点 $(a_n,b_n)$ 在 $y=2^x$ 的图象上得
\nopagebreak[4]
\[b_n=2^{a_n}>0 .\]
\Needspace{5\baselineskip}
因为 $\{a_n\}$ 是公差为 $d$ 的等差数列，$a_{n+1}-a_n=d$，所以
\nopagebreak[4]
\[\frac{b_{n+1}}{b_n}=\frac{2^{a_{n+1}}}{2^{a_n}}=2^{a_{n+1}-a_n}=2^d\]
是与 $n$ 无关的常数。故 $\{b_n\}$ 是首项为 $2^{a_1}$、公比为 $2^d$ 的等比数列。

\gksubpart{（2）}

\Needspace{5\baselineskip}
$f'(x)=2^x\ln2$，故曲线在点 $(a_2,b_2)=(a_2,2^{a_2})$ 处的切线为
\nopagebreak[4]
\[y-2^{a_2}=2^{a_2}\ln2\,(x-a_2).\]
\Needspace{5\baselineskip}
令 $y=0$，得切线在 $x$ 轴上的截距
\nopagebreak[4]
\[x=a_2-\frac1{\ln2} .\]
\Needspace{5\baselineskip}
由题设
\nopagebreak[4]
\[a_2-\frac1{\ln2}=2-\frac1{\ln2}\ \Longrightarrow\ a_2=2 .\]
\Needspace{5\baselineskip}
又 $a_1=1$，故公差 $d=a_2-a_1=1$，
\nopagebreak[4]
\[a_n=n,\qquad b_n=2^{a_n}=2^n .\]
\Needspace{5\baselineskip}
于是
\nopagebreak[4]
\[a_nb_n^2=n\cdot 4^n .\]

\Needspace{5\baselineskip}
记 $S_n=\displaystyle\sum_{k=1}^n k\cdot4^k$，则
\nopagebreak[4]
\[S_n=4+2\cdot4^2+3\cdot4^3+\cdots+n\cdot4^n,\]
\nopagebreak[4]
\[4S_n=4^2+2\cdot4^3+\cdots+(n-1)4^n+n\cdot4^{n+1}.\]
\Needspace{5\baselineskip}
两式相减：
\nopagebreak[4]
\[-3S_n=4+4^2+4^3+\cdots+4^n-n\cdot4^{n+1}
=\frac{4(4^n-1)}3-n\cdot4^{n+1},\]
\Needspace{5\baselineskip}
故
\nopagebreak[4]
\[S_n=\frac{n\cdot4^{n+1}-\frac{4(4^n-1)}3}{3}
=\frac{3n\cdot4^{n+1}-4^{n+1}+4}{9}
=\frac{(3n-1)4^{n+1}+4}{9}.\]

\gkanswer{答案（1）：证明。$b_n=2^{a_n}>0$，且 $\dfrac{b_{n+1}}{b_n}=2^{a_{n+1}-a_n}=2^d$ 为常数，故 $\{b_n\}$ 是公比为 $2^d$ 的等比数列。}

\gkanswer{答案（2）：$S_n=\dfrac{(3n-1)4^{n+1}+4}9$。切线截距 $a_2-\dfrac1{\ln2}=2-\dfrac1{\ln2}$ 给出 $a_2=2$，故 $a_n=n$，$b_n=2^n$，$a_nb_n^2=n\cdot4^n$，错位相减即得。}

\end{answer}
\begin{answer}
\gksolhead{20}{}

\textbf{题目}：

已知椭圆 $C:\dfrac{x^2}{a^2}+\dfrac{y^2}{b^2}=1$（$a>b>0$）的左焦点为 $F(-2,0)$，离心率为 $\dfrac{\sqrt6}3$。

（1） 求椭圆 $C$ 的标准方程；

（2） 设 $O$ 为坐标原点，$T$ 为直线 $x=-3$ 上一点，过 $F$ 作 $TF$ 的垂线交椭圆于 $P$，$Q$。当四边形 $OPTQ$ 是平行四边形时，求四边形 $OPTQ$ 的面积。

\Needspace{4\baselineskip}
\textbf{解析}

\gksubpart{（1）}

\Needspace{5\baselineskip}
由左焦点 $F(-2,0)$ 得 $c=2$。又离心率
\nopagebreak[4]
\[\frac ca=\frac{\sqrt6}3,\]
\Needspace{5\baselineskip}
故
\nopagebreak[4]
\[a=\frac{3c}{\sqrt6}=\frac{6}{\sqrt6}=\sqrt6,\qquad a^2=6,\]
\nopagebreak[4]
\[b^2=a^2-c^2=6-4=2 .\]
\Needspace{5\baselineskip}
椭圆 $C$ 的标准方程为
\nopagebreak[4]
\[\frac{x^2}6+\frac{y^2}2=1 .\]

\gksubpart{（2）}

\Needspace{5\baselineskip}
设 $T(-3,m)$，则 $\overrightarrow{FT}=(-1,m)$。与 $\overrightarrow{FT}$ 垂直的方向向量可取 $(m,1)$，因为
\nopagebreak[4]
\[(-1,m)\cdot(m,1)=-m+m=0 .\]
\Needspace{5\baselineskip}
过 $F(-2,0)$ 且垂直于 $TF$ 的直线为
\nopagebreak[4]
\[\begin{cases}x=-2+ms,\\y=s,\end{cases}\quad s\in\mathbf R .\]
\Needspace{5\baselineskip}
代入椭圆方程：
\nopagebreak[4]
\[\frac{(-2+ms)^2}6+\frac{s^2}2=1
\ \Longrightarrow\ (m^2s^2-4ms+4)+3s^2=6,\]
\Needspace{5\baselineskip}
即
\nopagebreak[4]
\[(m^2+3)s^2-4ms-2=0 .\]
\Needspace{5\baselineskip}
设 $P$，$Q$ 对应的参数为 $s_1$，$s_2$，则
\nopagebreak[4]
\[s_1+s_2=\frac{4m}{m^2+3},\qquad s_1s_2=-\frac2{m^2+3}.\]

\Needspace{5\baselineskip}
四边形 $OPTQ$ 是平行四边形，当且仅当对角线 $OT$ 与 $PQ$ 互相平分，即
\nopagebreak[4]
\[P+Q=O+T=T .\]
\Needspace{5\baselineskip}
比较横坐标：
\nopagebreak[4]
\[(-2+ms_1)+(-2+ms_2)=-3
\ \Longrightarrow\ m(s_1+s_2)=1
\ \Longrightarrow\ \frac{4m^2}{m^2+3}=1,\]
\Needspace{5\baselineskip}
故 $m^2=1$。比较纵坐标：
\nopagebreak[4]
\[s_1+s_2=m\ \Longrightarrow\ \frac{4m}{m^2+3}=m\ \Longrightarrow\ m(m^2-1)=0,\]
与 $m^2=1$ 一致，故 $m=\pm1$。

\Needspace{5\baselineskip}
此时
\nopagebreak[4]
\[S_{OPTQ}=\left|x_Py_Q-x_Qy_P\right|
=\left|(s_1-2)s_2-(s_2-2)s_1\right|
=2|s_1-s_2| .\]
\Needspace{5\baselineskip}
以 $m=1$ 为例，方程化为 $2s^2-2s-1=0$，故
\nopagebreak[4]
\[s_{1,2}=\frac{1\pm\sqrt3}2,\qquad |s_1-s_2|=\sqrt3,\]
\Needspace{5\baselineskip}
于是
\nopagebreak[4]
\[S_{OPTQ}=2\sqrt3 .\]
（$m=-1$ 时同理，方程化为 $2s^2+2s-1=0$，$|s_1-s_2|=\sqrt3$，面积同为 $2\sqrt3$。）

\gkanswer{答案（1）：$\dfrac{x^2}6+\dfrac{y^2}2=1$。$c=2$，由 $e=\dfrac{\sqrt6}3$ 得 $a=\sqrt6$，$b^2=a^2-c^2=2$。}

\gkanswer{答案（2）：$2\sqrt3$。设 $T(-3,m)$，过 $F$ 的垂线为 $x=-2+ms$，$y=s$，代入椭圆得 $(m^2+3)s^2-4ms-2=0$；由 $P+Q=T$ 得 $m^2=1$，面积 $=|x_Py_Q-x_Qy_P|=2|s_1-s_2|=2\sqrt3$。}

\end{answer}
\begin{answer}
\gksolhead{21}{}

\textbf{题目}：

已知函数 $f(x)=\mathrm e^x-ax^2-bx-1$，其中 $a$，$b\in\mathbf R$，$\mathrm e=2.71828\cdots$ 为自然对数的底数。

（1） 设 $g(x)$ 是函数 $f(x)$ 的导函数，求函数 $g(x)$ 在区间 $[0,1]$ 上的最小值；

（2） 若 $f(1)=0$，函数 $f(x)$ 在区间 $(0,1)$ 内有零点，证明：$\mathrm e-2<a<1$。

\Needspace{4\baselineskip}
\textbf{解析}

\gksubpart{（1）}

\nopagebreak[4]
\[f'(x)=\mathrm e^x-2ax-b,\]
\Needspace{5\baselineskip}
故
\nopagebreak[4]
\[g(x)=\mathrm e^x-2ax-b,\qquad g'(x)=\mathrm e^x-2a .\]
$g'$ 在 $[0,1]$ 上严格递增，分三种情形。

\Needspace{5\baselineskip}
若 $a\leqslant\dfrac12$，则 $2a\leqslant1\leqslant\mathrm e^x$，$g'(x)\geqslant0$，$g$ 在 $[0,1]$ 上递增，
\nopagebreak[4]
\[g_{\min}=g(0)=1-b .\]

\Needspace{5\baselineskip}
若 $a\geqslant\dfrac{\mathrm e}2$，则 $2a\geqslant\mathrm e\geqslant\mathrm e^x$，$g'(x)\leqslant0$，$g$ 在 $[0,1]$ 上递减，
\nopagebreak[4]
\[g_{\min}=g(1)=\mathrm e-2a-b .\]

\Needspace{5\baselineskip}
若 $\dfrac12<a<\dfrac{\mathrm e}2$，则 $g'(x)=0$ 有唯一解 $x_0=\ln(2a)\in(0,1)$，且 $g$ 在 $[0,x_0]$ 上递减、在 $[x_0,1]$ 上递增，故
\nopagebreak[4]
\[g_{\min}=g(\ln 2a)=2a-2a\ln(2a)-b .\]

\Needspace{5\baselineskip}
综上，
\nopagebreak[4]
\[g_{\min}=\begin{cases}
1-b,&a\leqslant\dfrac12,\\[4pt]
2a-2a\ln(2a)-b,&\dfrac12<a<\dfrac{\mathrm e}2,\\[4pt]
\mathrm e-2a-b,&a\geqslant\dfrac{\mathrm e}2 .
\end{cases}\]

\gksubpart{（2）}

\Needspace{5\baselineskip}
由 $f(1)=0$ 得
\nopagebreak[4]
\[\mathrm e-a-b-1=0\ \Longrightarrow\ b=\mathrm e-1-a,\]
\Needspace{5\baselineskip}
于是
\nopagebreak[4]
\[f(0)=1-0-0-1=0 .\]
\Needspace{5\baselineskip}
设 $x_0\in(0,1)$ 使 $f(x_0)=0$。则 $f$ 在 $[0,x_0]$ 与 $[x_0,1]$ 上均满足罗尔定理的条件，故存在
\nopagebreak[4]
\[\xi_1\in(0,x_0),\qquad \xi_2\in(x_0,1),\]
\Needspace{5\baselineskip}
使
\nopagebreak[4]
\[g(\xi_1)=f'(\xi_1)=0,\qquad g(\xi_2)=f'(\xi_2)=0,\]
且 $\xi_1<\xi_2$。

\Needspace{5\baselineskip}
又
\nopagebreak[4]
\[g''(x)=\mathrm e^x>0,\]
故 $g$ 在 $\mathbf R$ 上严格凸。

\Needspace{5\baselineskip}
设 $x<\xi_1$。把 $\xi_1$ 写成 $x$ 与 $\xi_2$ 的凸组合：取 $\lambda=\dfrac{\xi_2-\xi_1}{\xi_2-x}\in(0,1)$，则 $\xi_1=\lambda x+(1-\lambda)\xi_2$。由凸性
\nopagebreak[4]
\[g(\xi_1)\leqslant\lambda g(x)+(1-\lambda)g(\xi_2)=\lambda g(x),\]
\Needspace{5\baselineskip}
即 $0\leqslant\lambda g(x)$，故 $g(x)\geqslant0$。若 $g(x)=0$，则上式取等号，$g$ 在 $[x,\xi_2]$ 上是仿射函数，$g''=0$ 于 $(x,\xi_2)$，与 $g''=\mathrm e^x>0$ 矛盾。故
\nopagebreak[4]
\[g(x)>0\quad(x<\xi_1).\]
同法可证 $g(x)>0$ 对一切 $x>\xi_2$ 成立。

\Needspace{5\baselineskip}
因为 $0<\xi_1<\xi_2<1$，所以
\nopagebreak[4]
\[g(0)>0,\qquad g(1)>0 .\]
\Needspace{5\baselineskip}
而
\nopagebreak[4]
\[g(0)=1-b=1-(\mathrm e-1-a)=a+2-\mathrm e>0\ \Longrightarrow\ a>\mathrm e-2,\]
\nopagebreak[4]
\[g(1)=\mathrm e-2a-b=\mathrm e-2a-(\mathrm e-1-a)=1-a>0\ \Longrightarrow\ a<1 .\]
\Needspace{5\baselineskip}
故
\nopagebreak[4]
\[\mathrm e-2<a<1 .\]

\gkanswer{答案（1）：当 $a\leqslant\dfrac12$ 时 $g_{\min}=1-b$；当 $\dfrac12<a<\dfrac{\mathrm e}2$ 时 $g_{\min}=2a-2a\ln(2a)-b$；当 $a\geqslant\dfrac{\mathrm e}2$ 时 $g_{\min}=\mathrm e-2a-b$。由 $g'(x)=\mathrm e^x-2a$ 的零点位置分三种情形求得。}

\gkanswer{答案（2）：证明。由 $f(1)=0$ 得 $b=\mathrm e-1-a$，且 $f(0)=0$；对 $[0,x_0]$ 与 $[x_0,1]$ 用罗尔定理知 $g=f'$ 在 $(0,1)$ 内有两个零点，而 $g''=\mathrm e^x>0$ 使 $g$ 严格凸，故 $g(0)>0$、$g(1)>0$，即 $a+2-\mathrm e>0$、$1-a>0$，从而 $\mathrm e-2<a<1$。}

\end{answer}